\documentclass[11pt,a4paper]{article}

\usepackage[utf8]{inputenc}
\usepackage[T1]{fontenc}
\usepackage[spanish]{babel}
\usepackage[a4paper,margin=2.4cm]{geometry}
\usepackage{amsmath,amssymb}
\usepackage{bm}
\usepackage{booktabs}
\usepackage{enumitem}
\usepackage{xcolor}
\usepackage{hyperref}
\usepackage{tcolorbox}

\hypersetup{colorlinks=true, linkcolor=blue!50!black, urlcolor=blue!60!black}

\newcommand{\BHF}{B_{\mathrm{hf}}}
\newcommand{\dEQ}{\Delta E_{Q}}
\newcommand{\code}[1]{\texttt{\small #1}}
\DeclareMathOperator*{\argmin}{arg\,min}

\newtcolorbox{nota}[1][Nota]{colback=blue!4, colframe=blue!45!black,
  fonttitle=\bfseries, title={#1}, left=4pt, right=4pt, top=3pt, bottom=3pt}
\newtcolorbox{aviso}[1][Cuándo usarla]{colback=orange!7, colframe=orange!70!black,
  fonttitle=\bfseries, title={#1}, left=4pt, right=4pt, top=3pt, bottom=3pt}

\title{\textbf{Corrección por espesor en \emph{Mossbauer}}\\[3pt]
\large Modelo de absorbente grueso: beneficios, aplicabilidad y uso\\
\normalsize\textcolor{black!60}{(rama experimental \texttt{claude/experimentos-espesor})}}
\author{Departamento de Química Física · UAM}
\date{\today}

\begin{document}
\maketitle

\begin{abstract}\noindent
El modelo de ajuste estándar trata la muestra como un absorbente \emph{delgado}:
la transmisión es lineal en la absorción de cada componente. En muestras gruesas
o muy absorbentes esa aproximación falla —las líneas profundas se \emph{saturan}—
y sesga profundidades, anchuras y, sobre todo, las áreas (fracciones de fase).
La rama experimental añade una \emph{corrección por espesor} mediante un modelo de
saturación exponencial con un único parámetro físico $C$. Este documento explica
qué beneficios aporta, cuándo conviene aplicarla, cómo se aplica (qué cambia en el
ajuste y en la representación) y cómo se comporta con varios subespectros.
\end{abstract}

\tableofcontents
\bigskip

%======================================================================
\section{El problema: absorbente delgado vs. grueso}

El modelo histórico es lineal en la absorción:
\begin{equation}
  T_{\text{delgado}}(v) \;=\; b + s\,v \;-\; A_{\text{tot}}(v),
  \qquad A_{\text{tot}}(v)=\sum_{c} A_c(v),
  \label{eq:thin}
\end{equation}
con $b$ la línea base, $s$ la pendiente y $A_c$ la absorción de cada componente
(sextete/doblete/singlete). Es exacto sólo en el límite de absorbente
infinitamente fino. En una muestra real de espesor finito, la fracción de fotones
resonantes que llegan al detector decae \emph{exponencialmente} con la cantidad de
material resonante: las líneas más profundas absorben proporcionalmente menos de
lo que predice \eqref{eq:thin} (se \emph{saturan}) y, además, se \emph{ensanchan}.

\begin{nota}
El efecto de saturación de espesor no es un artefacto numérico: es física del
experimento de transmisión. Ignorarlo en muestras gruesas hace que el ajuste
``compense'' falsamente subiendo anchuras y desequilibrando áreas.
\end{nota}

%======================================================================
\section{Modelo de saturación exponencial}

La corrección sustituye \eqref{eq:thin} por
\begin{equation}
  \boxed{\;
  T(v) \;=\; b + s\,v \;-\; C\,\Bigl(1 - e^{-A_{\text{tot}}(v)/C}\Bigr)
  \;}
  \label{eq:thick}
\end{equation}
donde $C>0$ es la \emph{escala de saturación} (amplitud de absorción máxima
alcanzable, $\approx f_s\,b$, con $f_s$ la fracción libre de retroceso efectiva
de la fuente). Sus dos límites lo hacen interpretable y compatible hacia atrás:
\begin{itemize}[leftmargin=*]
  \item \textbf{Límite delgado} $C\to\infty$: como $C(1-e^{-x/C})\to x$, se
        recupera exactamente \eqref{eq:thin}. Es decir, el modelo delgado es el
        caso particular de espesor nulo.
  \item \textbf{Saturación} $C$ finito: la absorción efectiva
        $A_{\text{eff}}=C(1-e^{-A_{\text{tot}}/C})$ está acotada por $C$ y
        comprime las líneas profundas tanto más cuanto mayor es
        $A_{\text{tot}}/C$.
\end{itemize}
Implementado en \code{core/physics.py} (\code{total\_model}, parámetro
\code{sat\_scale}=$C$).

%======================================================================
\section{Beneficios}

\begin{enumerate}[leftmargin=*]
  \item \textbf{Áreas y fracciones de fase insesgadas.} La saturación suprime más
        la absorción de las fases mayoritarias (líneas profundas) que la de las
        minoritarias. Sin corregir, el cociente de áreas está sesgado a favor de
        las fases débiles. Al \emph{des-saturar}, las áreas recuperadas
        corresponden a las del absorbente delgado equivalente, que son las
        proporcionales a las abundancias (a igual $f$ de Lamb–Mössbauer).
  \item \textbf{Profundidades correctas.} El parámetro de profundidad deja de
        leer la absorción aparente (recortada) para representar la absorción real.
  \item \textbf{Anchuras menos infladas.} El ensanchamiento por espesor ya no se
        absorbe artificialmente en $\gamma$; las anchuras ajustadas se acercan a
        las intrínsecas.
  \item \textbf{Mejor $\chi^2$ en espectros profundos.} Donde el modelo delgado
        deja estructura sistemática en el residuo (en los mínimos), la saturación
        la elimina sin añadir componentes espurios.
\end{enumerate}

\begin{nota}[Verificación]
Sobre un espectro sintético saturado con $A_{\max}/C=2$, el ajuste con corrección
de espesor reconstruyó el espectro un $\sim27\%$ mejor (RMS) que el modelo
delgado, y recuperó la escala $C$ de partida ($C=0.41$ frente al verdadero
$0.40$). En saturación suave recuperó $\sim9\%$ más área (des-saturación).
\end{nota}

%======================================================================
\section{Cuándo aplicarla}

\begin{aviso}
La corrección de espesor \textbf{aporta} cuando la muestra es gruesa / muy
absorbente y \textbf{no es necesaria} (y no debe forzarse) cuando es fina.
\end{aviso}

Señales de que conviene activarla:
\begin{itemize}[leftmargin=*]
  \item Absorción profunda: el mínimo de transmisión baja bastante por debajo de
        la base (típicamente $\gtrsim 15$–$20\%$ de absorción), o espesor efectivo
        $t_a\gtrsim 1$.
  \item El ajuste delgado deja \emph{residuo estructurado en los mínimos} (el
        modelo no llega al fondo de las líneas profundas o las hace demasiado
        anchas).
  \item Se busca \emph{cuantificación de fases} fiable (cocientes de área), no
        sólo posiciones.
\end{itemize}
Cuándo \emph{no}:
\begin{itemize}[leftmargin=*]
  \item Espectros finos (absorción de pocos \%): el modelo lineal ya es exacto y
        $C$ queda indeterminado (tiende a valores grandes). Mantener modo delgado.
  \item Como diagnóstico: si al ajustar $C$ éste converge a un valor grande y el
        $\chi^2$ no mejora, la muestra es efectivamente delgada.
\end{itemize}

%======================================================================
\section{Cómo se aplica en el ajuste}

\subsection{Selección y parámetro}
En el panel de calibración hay un selector \emph{Absorbente: \code{thin} /
\code{thickness}} y un \emph{slider} de escala de saturación $C$
(\code{sat\_scale}), activo sólo en modo grueso. $C$ es un parámetro global
(como la base o la pendiente), ajustable o fijable.

\subsection{Modo discreto}
El optimizador no lineal (TRF) ya ajusta una función arbitraria, así que basta
con que el modelo directo use \eqref{eq:thick}: $C$ entra como un parámetro libre
más y se refina junto al resto. No cambia el algoritmo, sólo el forward.

\subsection{Modo distribución \texorpdfstring{$P(\BHF)$}{P(BHF)}}
Aquí la exponencial \emph{rompe la linealidad} en los pesos $P$, sobre la que se
apoya todo el solver regularizado. La solución es una \emph{transformada inversa}
que recupera la linealidad: dado el dato $y$, se invierte la saturación para
obtener la absorción lineal observada
\begin{equation}
  A_{\text{obs}}(v) \;=\; -\,C\,\ln\!\Bigl(1 - \tfrac{b + s\,v - y(v)}{C}\Bigr),
  \label{eq:invtransform}
\end{equation}
que \emph{sí} es lineal en $P$ ($A_{\text{obs}}=K\,P+\dots$). Se resuelve el
sistema lineal regularizado habitual (Tikhonov/TV, no negatividad) sobre
$A_{\text{obs}}$, y el modelo se \emph{re-satura} con \eqref{eq:thick} para
calcular residuos y la curva mostrada. El ruido se propaga por la transformada,
\begin{equation}
  \sigma_{A} \;=\; \sigma_y\,\frac{dA_{\text{obs}}}{dy} \;=\; \sigma_y\,e^{A_{\text{obs}}/C},
\end{equation}
de modo que los canales saturados pesan correctamente. Un lazo externo no lineal
(esquema separable tipo VARPRO) refina $(b,s,C)$ mientras el solve lineal interno
recupera $P$ en cada iteración.

\begin{nota}
La transformada \eqref{eq:invtransform} exige $b+s\,v-y < C$ en todo canal (la
absorción aparente no puede exceder la plateau de saturación); el programa recorta
el argumento al dominio válido $(0,1]$ antes del logaritmo.
\end{nota}

%======================================================================
\section{Qué cambia en la vista}

\begin{itemize}[leftmargin=*]
  \item La \textbf{curva de modelo} representada es la transmisión \emph{saturada}
        \eqref{eq:thick}: es la que se compara con los datos y la que genera el
        residuo.
  \item Las \textbf{profundidades/áreas} mostradas y exportadas son las del
        absorbente delgado equivalente (des-saturadas): por eso, al activar la
        corrección, las áreas relativas pueden cambiar respecto al ajuste delgado
        —es justamente la corrección buscada.
  \item Aparece el valor ajustado de $C$ (escala de saturación) entre los
        parámetros globales. Un $C$ grande indica muestra prácticamente delgada.
  \item En modo distribución, $P(\BHF)$ se muestra ya des-saturada; su área
        integрада refleja la fracción real.
\end{itemize}

%======================================================================
\section{Varios subespectros}

La saturación es una propiedad del \emph{absorbente completo}, no de cada fase por
separado: los fotones de una velocidad dada atraviesan todo el material y ``ven''
la suma de todas las absorciones a esa velocidad. Por eso \eqref{eq:thick} satura
la \textbf{absorción total} $A_{\text{tot}}=\sum_c A_c$ con \emph{una sola} escala
$C$ común a todos los subespectros, y no cada componente por su cuenta:
\begin{equation}
  T(v) = b + s\,v - C\Bigl(1 - \exp\!\bigl[-\tfrac{1}{C}\textstyle\sum_c A_c(v)\bigr]\Bigr).
\end{equation}
Consecuencias prácticas:
\begin{itemize}[leftmargin=*]
  \item \textbf{Acoplamiento físico correcto.} Donde dos subespectros solapan, la
        saturación es mayor (la absorción local es la suma); el modelo lo captura
        automáticamente.
  \item \textbf{Cuantificación de fases.} Tras des-saturar, los cocientes de área
        entre subespectros recuperan su valor insesgado, que es el objetivo en
        mezclas de fases (p. ej. varios entornos de Fe).
  \item \textbf{Un único $C$.} No se ajusta una escala por componente (sería
        no físico y degeneraría con las profundidades): $C$ es global y describe el
        espesor efectivo de la muestra.
  \item \textbf{Compatibilidad.} En modo distribución, los sextetes ``nítidos''
        añadidos y la propia distribución comparten la misma corrección, porque
        todos contribuyen a $A_{\text{tot}}$.
\end{itemize}

\begin{nota}[Limitación del modelo]
Es una saturación exponencial efectiva (Lambert–Beer), que captura el efecto
dominante del espesor (compresión, ensanchamiento, áreas insesgadas) con un único
parámetro. No es la integral de transmisión completa de Margulies–Ehrman
(convolución con la línea de la fuente), que sería un refinamiento de segundo
orden raramente necesario en la práctica.
\end{nota}

%======================================================================
\section{Resumen}

\begin{center}
\begin{tabular}{@{}ll@{}}
\toprule
\textbf{Aspecto} & \textbf{Corrección por espesor}\\
\midrule
Modelo & $T=b+s v - C(1-e^{-A_{\text{tot}}/C})$ \\
Parámetro nuevo & $C$ (escala de saturación), global, $C\to\infty$ = delgado\\
Cuándo & muestras gruesas / absorción profunda / cuantificación de fases\\
Discreto & $C$ libre en el TRF; sólo cambia el modelo directo\\
Distribución & transformada inversa $A_{\text{obs}}=-C\ln(1-(b+sv-y)/C)$ + VARPRO $(b,s,C)$\\
Vista & curva saturada; áreas/$P$ des-saturadas; se muestra $C$\\
Subespectros & una sola $C$ satura la suma $\sum_c A_c$ (acoplamiento físico)\\
Beneficio clave & áreas/fracciones de fase insesgadas, anchuras y $\chi^2$ mejores\\
\bottomrule
\end{tabular}
\end{center}

\end{document}
